\section{Corrected coordinate grids}\label{sec:grids}

This section shows how a finite dynamic program yields a valid lower bound
for the continuous mixed-integer problem, and how its coordinate
min-marginals remove parts of the domain without losing any global minimizer.
Only Definition~\ref{def:curvature} is used: no convexity, uniqueness or
growth.

\subsection{Corrected grids and cell bounds}
A \emph{subbox} is a product $B=\prod_iB_i\subseteq X$ in which
$B_i=[\ell_i',u_i']$ or $[\ell_i',u_i']\cap\Z$ has the type of $X_i$, with
integer endpoints for $i\in I_Z$. A \emph{grid} for $B$ is a product
$G=\prod_iG_i$ of finite sets $G_i\subseteq B_i$ that contain $\ell_i'$ and
$u_i'$. Consecutive nodes $a<a'$ of $G_i$ form a \emph{grid interval}
$[a,a']$. Its \emph{effective width} is $\Delta_i([a,a'])=0$ if $i\in I_Z$ and
$a'-a=1$, and $a'-a$ otherwise; an integer interval of length one has no
feasible point in its interior.

\begin{definition}[Corrections, bound and min-marginals]\label{def:corr}
For $v\in G_i$ let $w_i(v)$ be the largest effective width of a grid interval
with endpoint $v$ ($w_i(v)=0$ if there is none). Put
\begin{equation}\label{eq:correction}
d_i(v)=\frac{L_i}{8}w_i(v)^2,\qquad D(y)=\sum_{i=1}^nd_i(y_i),\qquad
Q(y)=F(y)-D(y)\qquad(y\in G),
\end{equation}
\begin{equation}\label{eq:minmarginal}
\beta=\min_{y\in G}Q(y),\qquad m_i(v)=\min\{Q(y):y\in G,\ y_i=v\}\quad(v\in G_i).
\end{equation}
The function $m_i$ is the \emph{min-marginal} of $Q$ in coordinate $i$.
\end{definition}

The corrections are unary, so $Q$ has the same scopes as $F$ and the same
tree decomposition applies to it. A \emph{cell} of $G$ is a product
$C=\prod_iI_i$ of grid intervals $I_i$ of $G_i$; its \emph{cell bound} is
\[
\beta(C)=\min_{v\in\mathrm{vert}(C)}F(v)-\sum_i\frac{L_i}{8}\Delta_i(I_i)^2,
\qquad\mathrm{vert}(C)=\prod_i\{\text{endpoints of }I_i\}.
\]
Subtracting $\frac{L_i}2(x_i-\mathrm{mid}\,I_i)^2$ in every coordinate makes
$F$ concave along each coordinate of $C$, so its minimum over $C$ is attained
at a vertex, and $\beta(C)$ is a lower bound for $F$ on $C$. For a single box
this is the vertex bound of Bajaj and Hasan \cite[Theorem~1]{BajajHasan2020},
an instance of edge-concave underestimation
\cite{Tardella2004,MeyerFloudas2005,Hasan2018}; the term
$\frac{L_i}8\Delta_i^2$ is the maximum separation of an $\alpha$BB
underestimator with $\alpha_i=L_i/2$ \cite{MaranasFloudas1994,AdjimanEtAl1998},
used with the opposite sign. The next proposition shows that the unary
correction~\eqref{eq:correction} computes the best of these bounds over all
cells of a product grid at once.

\begin{proposition}[Cellwise form of the corrected grid]\label{prop:cellwise}
Let $G$ be a grid for a subbox $B$.
\begin{enumerate}[label=(\alph*)]
\item For every cell $C$ and every $x\in C\cap X$, $F(x)\ge\beta(C)$. Every
$x\in B$ lies in some cell.
\item $\beta=\min_C\beta(C)$, the minimum over all cells. In particular
$\beta\le\min_BF$.
\item For a grid interval $I=[a,a']$ of coordinate $i$, put
$\beta_i(I)=\min\{\beta(C):C_i=I\}$. Then
\[
\beta_i(I)=\min_{v\in\{a,a'\}}\bigl(m_i(v)+d_i(v)\bigr)-\frac{L_i}8\Delta_i(I)^2
\ \ge\ \min\{m_i(a),m_i(a')\},
\]
and $F(x)\ge\beta_i(I)$ for every $x\in B$ with $x_i\in I$.
\end{enumerate}
\end{proposition}

\begin{proof}
(a) Let $C=\prod_iI_i$ with $I_i=[a_i,a_i']$, let $J=\{i:\Delta_i(I_i)>0\}$,
$c_i=(a_i+a_i')/2$, and
$\varphi(y)=F(y)-\sum_{i\in J}\frac{L_i}{2}(y_i-c_i)^2\le F(y)$. If $i\notin J$,
then $i\in I_Z$ and $a_i'=a_i+1$, so $x_i\in\{a_i,a_i'\}$. If $i\in J$, then
for fixed values of the other coordinates,
$t\mapsto\varphi(\dots,t,\dots)=\bigl(F(\dots,t,\dots)-\frac{L_i}2t^2\bigr)
+L_ic_it+\mathrm{const}$ is concave on $I_i$ by
Definition~\ref{def:curvature}, so its minimum over $I_i$ is attained at an
endpoint. Replacing the coordinates in $J$ one at a time by such endpoints
keeps the point in $C\cap X$ and does not increase $\varphi$, and ends at a
vertex $v$ with $\varphi(v)\le\varphi(x)$. Since $(v_i-c_i)^2=\Delta_i(I_i)^2/4$
for $i\in J$,
$F(x)\ge\varphi(x)\ge\varphi(v)=F(v)-\sum_i\frac{L_i}8\Delta_i(I_i)^2\ge\beta(C)$.
Every $x_i\in B_i$ lies in a grid interval because $G_i$ contains both endpoints
of $B_i$.

(b) For $v\in G_i$ let $\mathcal I_i(v)$ be the set of grid intervals with
endpoint $v$. The pairs $(C,y)$ with $y\in\mathrm{vert}(C)$ are exactly the
grid points $y$ together with an independent choice of $I_i\in\mathcal I_i(y_i)$
for every $i$. Hence
\[
\min_C\beta(C)=\min_{y\in G}\Bigl[F(y)-\sum_i\max_{I\in\mathcal I_i(y_i)}
\frac{L_i}8\Delta_i(I)^2\Bigr]=\min_{y\in G}\bigl[F(y)-D(y)\bigr]=\beta .
\]
(c) The same computation with the $i$th interval fixed to $I$ gives the
formula, because $Q(y)=F(y)-\sum_{k\ne i}d_k(y_k)-d_i(v)$ when $y_i=v$. Since
$I\in\mathcal I_i(v)$ for $v\in\{a,a'\}$, $d_i(v)\ge L_i\Delta_i(I)^2/8$, which
gives the inequality. The last claim follows from (a) with $C_i=I$.
\end{proof}

Enumerating cells costs $\prod_i(\abs{G_i}-1)$ cells with up to $2^n$ vertex
evaluations each, as in box-by-box branch-and-bound. Because $D$ is a sum of
unary terms, the same number is the minimum of a factorized function on the
product grid, and it is computed by one min-sum computation on the given tree
decomposition with tables of size $\max_i\abs{G_i}^p$. The same bound also
follows from randomized rounding: rounding each coordinate of $x$
independently to the endpoints of its grid interval, with probabilities that
preserve the mean, gives a random grid point $Y$ with $\E Q(Y)\le F(x)$
(see the proof of Lemma~\ref{lem:tu-round} for the correlated version used with
constraints).

\subsection{Dynamic programming and min-marginals}
Let $q_t$ be the sum of the factors assigned to bag $t$ minus the corrections
$d_i$ of the coordinates whose home bag is $t$, regarded as a table on
$G_{B_t}=\prod_{i\in B_t}G_i$; thus $Q=\sum_tq_t$ on $G$. For adjacent nodes
$t,u$ let $S_{tu}=B_t\cap B_u$ and define messages
\begin{equation}\label{eq:messages}
M_{t\to u}(y_{S_{tu}})=\min_{y_{B_t\setminus S_{tu}}}\Bigl[q_t(y_{B_t})
+\sum_{v\in\mathcal N(t)\setminus\{u\}}M_{v\to t}(y_{S_{vt}})\Bigr],
\end{equation}
where $\mathcal N(t)$ is the set of neighbours of $t$ and an empty separator
carries a single scalar. Put
$A_t=q_t+\sum_{v\in\mathcal N(t)}M_{v\to t}$.

\begin{lemma}[Two-pass dynamic programming]\label{lem:dp}
For every bag $t$ and every $y_{B_t}\in G_{B_t}$,
$A_t(y_{B_t})=\min\{Q(z):z\in G,\ z_{B_t}=y_{B_t}\}$. Hence $\beta=\min A_t$, and
$m_i(v)$ is the minimum of $A_t$ over $y_i=v$ for any bag $t\ni i$. With
$K=\max_i\abs{G_i}$, all messages, $\beta$, a minimizer and all min-marginals
are obtained with $O\bigl((\abs{\mathcal A}+n+N)K^p\bigr)$ additions and
comparisons, $O(p)$ index work per table entry, and the evaluation of every
factor at the at most $K^{\abs{S_a}}$ grid points of its scope.
\end{lemma}

\begin{proof}
For an oriented edge $t\to u$ let $T_{t\to u}$ be the component of
$T\setminus\{tu\}$ containing $t$, and $V_{t\to u}$ the set of coordinates
occurring in bags of $T_{t\to u}$ but not in $S_{tu}$. By induction on
$\abs{T_{t\to u}}$, $M_{t\to u}(y_{S_{tu}})$ is the minimum of
$\sum_{r\in T_{t\to u}}q_r$ over $y_{V_{t\to u}}$: the trees $T_{v\to t}$,
$v\in\mathcal N(t)\setminus\{u\}$, partition $T_{t\to u}\setminus\{t\}$, and by
running intersection a coordinate occurring both in $T_{v\to t}$ and outside it
lies in $S_{vt}$, so the sets $V_{v\to t}$ are disjoint from each other and
from $B_t$ and the minimization splits. The same argument with all neighbours
of $t$ gives the formula for $A_t$. Messages are computed towards a root and
back; at bag $t$ the sum of all incoming messages is formed once, and each
outgoing message is obtained from $A_t-M_{u\to t}$ by one pass, with no
product over children. A minimizer is recovered from the root outwards.
\end{proof}

This is the classical nonserial dynamic programming, bucket elimination or
junction-tree computation \cite{BerteleBrioschi1972,Dechter1999,
LauritzenSpiegelhalter1988}, and two passes give all min-marginals
\cite{WainwrightJaakkolaWillsky2005,KohliTorr2006}. The only new element is
that its input is the corrected objective $Q$.

\subsection{Filtering and certificates}

\begin{definition}[Filtering step]\label{def:filter}
Let $G$ be a grid for a subbox $B$ and $U\in\R$. For $i\in P$, \emph{retain}
a grid interval $[a,a']$ of coordinate $i$ if
\begin{equation}\label{eq:retain}
\min\{m_i(a),m_i(a')\}\le U ,
\end{equation}
and let $B_i'$ be the smallest interval containing the retained intervals
(intersected with $\Z$ for $i\in I_Z$). For $i\notin P$ let $B_i'=B_i$. The
filtered subbox is $B'=\prod_iB_i'$.
\end{definition}

\begin{proposition}[Safe filtering]\label{prop:filter}
Every $x\in B\setminus B'$ satisfies $F(x)>U$. Hence, if $U\ge\min_BF$, then
$B'$ contains every $x\in B$ with $F(x)\le U$, in particular every minimizer
of $F$ over $B$; and every minimizer $y$ of $Q$ over $G$ lies in $B'$.
\end{proposition}

\begin{proof}
If $x\in B\setminus B'$, some $i\in P$ has $x_i\notin B_i'$, so every grid
interval containing $x_i$ was not retained, and
Proposition~\ref{prop:cellwise}(c) gives $F(x)\ge\min\{m_i(a),m_i(a')\}>U$
for such an interval. For a minimizer $y$ of $Q$,
$m_i(y_i)\le Q(y)=\beta\le\min_BF\le U$, and $y_i$ is an endpoint of a grid
interval of coordinate $i$, which is therefore retained.
\end{proof}

A method that filters repeatedly produces nested subboxes. The following
certificate records only what is needed to check the result.

\begin{definition}[Path certificate]\label{def:cert}
A \emph{path certificate} consists of grids $G^{(0)},\dots,G^{(k)}$, a number
$\beta$ and a point $\hat x$. Let $B^{(j)}$ be the subbox spanned by
$G^{(j)}$. It is \emph{valid} if $\hat x\in X$, $B^{(0)}=X$,
$B^{(j+1)}\subseteq B^{(j)}$, and, with $Q^{(j)}$ and $m^{(j)}$ computed from
$G^{(j)}$ and the curvature bounds,
\begin{enumerate}[label=(C\arabic*)]
\item for $j<k$, every grid interval $[a,a']$ of $G^{(j)}_i$ not contained in
$B^{(j+1)}_i$ has $\min\{m^{(j)}_i(a),m^{(j)}_i(a')\}\ge\beta$;
\item $\min_{G^{(k)}}Q^{(k)}\ge\beta$.
\end{enumerate}
\end{definition}

\begin{theorem}[Path certificates]\label{thm:certificate}
If a path certificate is valid, then $\beta\le\OPT\le F(\hat x)$. Its
validity is checked by $k+1$ dynamic programs (Lemma~\ref{lem:dp}) in exact
arithmetic, given the curvature bounds; for a quadratic the checker computes
$L_i=\max\{H_{ii},0\}$ itself. No growth or uniqueness assumption is involved.
\end{theorem}

\begin{proof}
Let $x\in X$ and let $j$ be the largest index with $x\in B^{(j)}$. If $j=k$,
Proposition~\ref{prop:cellwise}(b) and (C2) give $F(x)\ge\beta$. Otherwise some
$x_i\notin B^{(j+1)}_i$; a grid interval of $G_i^{(j)}$ containing $x_i$ is not
contained in $B_i^{(j+1)}$, and (C1) with Proposition~\ref{prop:cellwise}(c)
gives $F(x)\ge\beta$.
\end{proof}

If the subboxes are produced by the filtering step with thresholds $U_j$ that
are values of feasible points, and $\beta$ is the corrected minimum of the
last grid, then the record is a valid path certificate: every minimizer lies
in every $B^{(j)}$ by Proposition~\ref{prop:filter}, so $\beta\le\OPT\le U_j$,
and a removed interval has endpoint bound larger than $U_j\ge\beta$. Thus a
certificate consists of the stage grids (or the boxes with the centers,
meshes and grading that generate them), one bound and one feasible point.
Messages, min-marginal tables and intermediate incumbents need not be stored,
and stages that remove nothing may be dropped. The stages that remove points
are needed: the last grid bounds $F$ only on $B^{(k)}$. A checker repeats
the dynamic programs, so checking costs about as much as the successful run;
Section~\ref{sec:computation} reports both costs. The certificate is a
branch-and-bound proof whose tree is a path; its logic is that of checkable
branch-and-bound certificates \cite{CheungGleixnerSteffy2017}. Removing an
interval whose bound exceeds an incumbent is optimality-based bound
tightening \cite{RyooSahinidis1996,BelottiEtAl2009,GleixnerEtAl2017}; here all
interval bounds come from one pair of message passes.
